\documentclass[10pt,a4paper]{amsart}
\usepackage[margin=19mm]{geometry}
\usepackage{amsmath,amssymb,mathtools}
\usepackage[scr=rsfs,cal=euler]{mathalfa}
\usepackage{xfrac}
\usepackage[unicode,hidelinks,bookmarks=false]{hyperref}
\newcommand{\ie}{i.e.\ }
\newcommand{\cf}{cf.\ }
\newcommand{\resp}{respectively\ }
\newcommand{\Subsection}{\subsection}
\providecommand{\nobreakdash}{\nobreak\hbox{-}}
\setlength{\emergencystretch}{2em}
\multlinegap0pt
\usepackage{mathtools}
\usepackage[shortlabels]{enumitem}
\usepackage{tikz}
\newtheorem{proposition}{Proposition}
\usepackage{cleveref}
\crefname{proposition}{Proposition}{Propositions}
\newcommand{\EE}{\mathbb E}
\newcommand{\PP}{\mathbb P}
\newcommand{\pP}{\mathscr{P}}
\title{A Necessary and Sufficient Condition for Global Stability of Order-Preserving Semigroups and Maps}
\author{Takashi Kamihigashi, John Stachurski}
\date{Source: arXiv:2604.03979v1 (CC BY 4.0)}
\begin{document}
\maketitle

\noindent\emph{Source and licence.} A Necessary and Sufficient Condition for Global Stability of Order-Preserving Semigroups and Maps. Version arXiv:2604.03979v1 (CC BY 4.0), distributed by arXiv under the Creative Commons Attribution 4.0 International licence, http://creativecommons.org/licenses/by/4.0/, as recorded in the arXiv licence metadata for this exact version and shown on the abstract page https://arxiv.org/abs/2604.03979v1. Full licence notice: \texttt{licenses/arxiv-2604.03979-CC-BY-4.0.txt}.\par\smallskip

\noindent\textit{Editorial scope.} Counted unit: the article's proposition p:orwomnew (source0.tex lines 3116--3120). The article's complete original proof of it is delivered with it (source0.tex lines 3122--3221, one \texttt{proof} environment of 100 lines). No mathematics is written, repaired or re-derived: the statement and the proof are the article's own lines, in the article's own notation, and no retained line is substituted or re-flowed.  Retained uncounted as the source-local context the proof needs: lines 1711--1714.  Nothing was omitted from any retained span, so the package carries no omission record.  No retained line contains a citation, so the wrapper carries no bibliography and no bibliography entry.  No retained line contains a figure environment, an image-inclusion command or drawing code, so no image is delivered and the package declares no asset dependency.  Editorial formatting only, in addition: an AMS \texttt{amsart} preamble that restates the article's own declarations and macros used by the retained lines, namely the environments \texttt{proposition} on separate counters, together with the standard packages those lines need.  The retained environments print in this document as Proposition 1 in order of appearance, whereas the article prints the same items with the numbering of its own full document.  The complete native source of the article as distributed in the arXiv e-print for this exact version is bundled unchanged as \texttt{sources/197852/source0.tex}.
        \begin{equation}\label{eq:omb}
            \PP_{x,x'} \{ \tau < \infty\} = 1
            \quad \text{for all } x, x' \in S.
        \end{equation}

\begin{proposition}\label{p:orwomnew}
    Let $(P_t)$ be an increasing transition probability function on $S$.
    If $(P_t)$ is weakly order mixing, then $(P_t)$ is
    asymptotically contractive on $\pP(S)$ with respect to the metric $\beta$.
\end{proposition}

\begin{proof}[Proof of \cref{p:orwomnew}]
    Fix $x, x' \in S$ and $h \in ibS$ with $|h| \leq 1$.  Let $((X_t),
    (X_t'))$ be as in the definition of weakly order mixing and $\tau$ be as
    defined in \eqref{eq:omb}. We have
    %
    \begin{equation*}
        \delta_x P_t h - \delta_{x'} P_t h
        = \EE_x  h(X_t) - \EE_{x'} h(X_t')
        = \EE_x \EE_\tau h(X_t) - \EE_{x'} \EE_\tau h(X_t')
    \end{equation*}
    %
    By construction, $(X_t)$ and $(X_t')$ are defined on the same probability
    space, so we can rewrite this as
    %
    \begin{equation*}
        \delta_x P_t h - \delta_{x'} P_t h
        = \EE_{x,x'} [\EE_\tau h(X_t) - \EE_\tau h(X_t')]
        = A + B,
    \end{equation*}
    %
    where
    %
    \begin{align*}
        A & \coloneq \EE_{x,x'} [\EE_\tau h(X_t) - \EE_\tau h(X_t')] \{\tau \leq t\}
        \\
        B & \coloneq \EE_{x,x'} [\EE_\tau h(X_t) - \EE_\tau h(X_t')] \{\tau > t\}.
    \end{align*}
    %
    Using the strong Markov property and the fact that $\tau$ is a finite
    stopping time, we can write $A$ as
    %
    \begin{equation*}
        A = \EE_{x,x'} [(P_{t-\tau}h)(X_\tau) - (P_{t-\tau} h)(X_\tau')] \{\tau \leq t\}
    \end{equation*}
    %
    Since $h$ is increasing, $X_\tau \leq X_\tau'$ and $(P_t)$ is increasing, we
    must have $A \leq 0$.  Moreover,
    %
    \begin{equation*}
        B \leq \EE_{x,x'} |\EE_\tau h(X_t) - \EE_\tau h(X_t')| \{\tau > t\}
        \leq 2 \PP_{x,x'} \{\tau > t\}.
    \end{equation*}
    %
    As a result,
    %
    \begin{equation*}
        \delta_x P_t h - \delta_{x'} P_t h \leq 2 \PP_{x,x'} \{\tau > t\}.
    \end{equation*}
    %
    Since the one-sided bound holds for all $x, x' \in S$, applying it
    to the pair $(x', x)$ gives $\delta_{x'} P_t h - \delta_x P_t h
    \leq 2 \PP_{x',x} \{\tau > t\}$, and hence
    %
    \begin{equation*}
        |\delta_x P_t h - \delta_{x'} P_t h| \leq 2 \max(\PP_{x,x'} \{\tau > t\},\, \PP_{x',x} \{\tau > t\}).
    \end{equation*}
    %
    Since $h \in ibS$ with $|h| \leq 1$ was arbitrary, we can take the supremum to get
    %
    \begin{equation}\label{eq:acd}
        \beta(\delta_x P_t , \delta_{x'} P_t) \leq 2 \max(\PP_{x,x'} \{\tau > t\},\, \PP_{x',x} \{\tau > t\})
        \quad \text{for all } x,x' \in S.
    \end{equation}
    %
    We wish to extend this to arbitrary (nondegenerate) initial conditions.
    To this end, fix $\phi, \psi \in \pP(S)$. Let $(X, X')$ be any coupling with marginals
    $\phi$ and $\psi$. For $h \in ibS$ with $|h| \leq 1$, we have
    %
    \begin{equation*}
        (\phi P_t)(h) - (\psi P_t)(h)
        = \EE[(P_t h)(X)] - \EE[(P_t h)(X')]
        = \EE[(\delta_X P_t)(h) - (\delta_{X'} P_t)(h)].
    \end{equation*}
    %
    Taking absolute values and using \eqref{eq:acd}, we get
    %
    \begin{equation*}
        |(\phi P_t)(h) - (\psi P_t)(h)| \leq \EE[\beta(\delta_X P_t, \delta_{X'} P_t)].
    \end{equation*}
    %
    The right-hand side does not depend on $h$, so taking the supremum over $h \in ibS$ with $|h| \leq 1$ gives
    %
    \begin{equation*}
        \beta(\phi P_t, \psi P_t) \leq \EE[\beta(\delta_X P_t, \delta_{X'} P_t)].
    \end{equation*}
    %
    By \eqref{eq:acd},
    %
    \begin{equation}\label{eq:wombd}
        \beta(\phi P_t, \psi P_t) \leq 2\EE[\max(\PP_{X,X'}\{\tau > t\},\, \PP_{X',X}\{\tau > t\})].
    \end{equation}
    %
    Since $\PP_{x,x'}\{\tau < \infty\} = 1$ by \eqref{eq:omb} for all $x, x' \in S$,
    both $\PP_{x,x'}\{\tau > t\}$ and $\PP_{x',x}\{\tau > t\}$ converge to $0$
    as $t \to \infty$ for each $(x,x')$.
    By the dominated convergence theorem, the right-hand side of
    \eqref{eq:wombd} converges to $0$.
    Since $\phi, \psi \in \pP(S)$ were arbitrary, $(P_t)$ is asymptotically
    contractive.
\end{proof}

\end{document}
